

\documentclass[12pt]{article}


% ---------- Font (requires XeLaTeX or LuaLaTeX compiler) ----------
\usepackage{fontspec}
\setmainfont{Montserrat}   % or change to Roboto
% \setmainfont{Roboto}     % <- uncomment this line (and comment the one above) to use Roboto instead

% ---------- Packages ----------
\usepackage[a4paper, margin=1in]{geometry}
\usepackage{graphicx}
\usepackage{xcolor}
\usepackage{enumitem}
\usepackage{titlesec}
\usepackage{parskip}
\usepackage{hyperref}
\usepackage{fancyhdr}

% ---------- Packages ----------
\usepackage[a4paper, margin=1in]{geometry}
\usepackage{graphicx}
\usepackage{xcolor}
\usepackage{enumitem}
\usepackage{titlesec}
\usepackage{parskip}
\usepackage{hyperref}
\usepackage{fancyhdr}

% ---------- Colors matching the annotation arrows ----------
\definecolor{arrowblue}{RGB}{31, 86, 214}
\definecolor{arrowred}{RGB}{214, 40, 40}
\definecolor{arrowgreen}{RGB}{27, 140, 60}
\definecolor{arrowpurple}{RGB}{128, 0, 200}
\definecolor{arroworange}{RGB}{255, 140, 0}
\definecolor{arrowgold}{RGB}{170, 130, 0}
\definecolor{navy}{RGB}{26, 26, 46}

% ---------- Section title style ----------
\titleformat{\section}{\large\bfseries\color{navy}}{\thesection}{1em}{}
\titlespacing*{\section}{0pt}{12pt}{8pt}

\pagestyle{fancy}
\fancyhf{}
\fancyfoot[C]{\thepage}
\renewcommand{\headrulewidth}{0pt}

\begin{document}

% ================= COVER PAGE =================
\begin{titlepage}
\centering
\vspace*{2cm}

{\large University of Sindh -- IMCS}\\[1.5cm]

{\Huge \bfseries \color{navy} HCI and Computer Graphics}\\[0.5cm]
{\Large Assignment \#1}\\[1cm]

\rule{0.7\textwidth}{0.8pt}\\[1.5cm]

\begin{tabular}{ll}
\vspace{12pt}
\textbf{Name:} & Ahtsham Nizamani \\[0.3cm] \vspace{12pt}
\textbf{Roll No:} & 2K24/CSE/23 \\[0.3cm] \vspace{12pt}
\textbf{Semester:} & 6th \\[0.3cm] \vspace{12pt}
\textbf{Subject:} & HCI and Computer Graphics \\[0.3cm] \vspace{12pt}
\textbf{Submitted To:} & Sir Rajesh \\
\end{tabular}

\vfill
\end{titlepage}

% ================= PART 1: HCI FOCUS =================
\begin{center}\section*{HCI - Assignment \#1: 3D Map App Overview --- Part \#1}\end{center}

This assignment part focuses on an interface analysis of Google Earth, as 3D Map Overview.\vspace{12pt}

\begin{center}
\includegraphics[width=0.95\textwidth]{annotated_image.png} \vspace{12pt}
\end{center}

\begin{itemize}[leftmargin=1.2em, itemsep=6pt]
    \item \textbf{\textcolor{arrowblue}{Blue Arrow (search bar):}} The \textcolor{arrowblue}{blue} arrow points to Google Earth's in-app search field --- the actual input point where a user types a location to search.

    \item \textbf{\textcolor{arrowred}{Red Arrow (URL bar):}} The \textcolor{arrowred}{red} arrow points from the map view up to the browser's address bar, showing how the URL updates live to reflect the current search (``Rajo Nizamani'') --- a form of system feedback.

    \item \textbf{\textcolor{arrowgreen}{Zoom Controls (Green Arrow):}} The \textcolor{arrowgreen}{green} arrow points to the bottom-right corner of the interface, highlighting the ``+'' (zoom in) and ``-'' (zoom out) controls. These allow users to adjust the map scale for detailed inspection or broader geographical context.
\end{itemize}

\vspace{12pt}\textbf{This is the first part of the assignment.}

\newpage

% ================= PART 2: CG FOCUS =================
\begin{center}\section*{HCI - Assignment \#1: 3D Map App --- Part 2 (CG Focus)}\end{center}

This part examines the computer graphics techniques of Google Earth's Street View. \vspace{12pt}

\begin{center}
\includegraphics[width=0.95\textwidth]{HCI_Assignment1_Part2_CG_Annotated (1).png}
\end{center}

\begin{itemize}[leftmargin=1.2em, itemsep=6pt] 
    \vspace{12pt} \item \textbf{\textcolor{arrowpurple}{Perspective Camera View (Purple Arrow):}} The \textcolor{arrowpurple}{purple} arrow points down S Homan Ave to where the road and buildings converge, showing perspective projection --- distant objects shrink, creating the illusion of depth.

    \item \textbf{\textcolor{arroworange}{Texture Mapping (Orange Arrow):}} The \textcolor{arroworange}{orange} arrow points to the plain siding on the beige building's exterior, where photographic material imagery is mapped onto its 3D surface. This gives the flat building geometry a realistic, detailed appearance.

    \item \textbf{\textcolor{arrowgold}{Light Source \& Shadows (Gold Arrow):}} The \textcolor{arrowgold}{gold} arrow points to the shadow cast on the sidewalk by nearby trees/structures, implying a single light source (the sun) and helping ground objects in 3D space.
\end{itemize}

\vspace{12pt} \textbf{This is the second part of the assignment.}

\newpage

% ================= PART 3: REFLECTION =================
\begin{center}\section*{HCI - Assignment \#1: Reflection Questions --- Part \#3}\end{center} \vspace{12pt}

\textbf{Question 1:} How does the Computer Graphics quality (e.g., realistic lighting or 3D view) help or hinder the user's ability to interact with the software?

\textbf{Answer:} The realistic satellite imagery and 3D view in Google Earth actually make it a lot easier to recognize real-world places and get a feel for how everything is laid out, so navigating feels pretty natural. That said, the tilted 3D angle isn't always helpful --- it can make text labels harder to read and distances tricky to judge, since the perspective kind of warps the scale of things.

\vspace{0.5cm}

\textbf{Question 2:} What is one simple HCI change you would make to improve the user interface?

\textbf{Answer:} I'd make the zoom controls more obvious --- right now they're just small icon buttons that are easy to overlook. Adding a visible ``+/-'' label next to them would make it clearer at a glance what they do, especially for someone using the app for the first time.

\vspace{1cm}
\textbf{This is the third part of the assignment.}

\end{document}
